\section{Analytic types and restriction}\label{sec:analytic-types}\index{Lelong number}\index{Analytic singularity}

The simplest singularities of quasi-plurisubharmonic functions are logarithmic poles along the common zeros of finitely many holomorphic functions. The logarithmic growth of a general singularity at a point is measured by its Lelong number. Let $u$ be a quasi-plurisubharmonic function on a complex manifold and let $x$ be a point. Choose holomorphic coordinates centered at $x$ and a smooth function $q$ such that $p=u+q$ is plurisubharmonic near $x$, and write $B(x,r)$ for the Euclidean coordinate ball of radius $r$. The \emph{Lelong number} of $u$ at $x$ is
\[
 \nu(u,x)=\lim_{r\downarrow0}
       \frac{\sup_{B(x,r)}p}{\log r^2}.
\]
By the local growth theorem recalled in \cite{SGPT}, the limit exists and does not depend on the coordinates or on the correction $q$. With this normalization, $\log|z|^2$ has Lelong number one at the origin. For an irreducible analytic subset $Z$, the \emph{generic Lelong number} of $u$ along $Z$ is
\[
 \nu(u,Z)=\inf_{x\in Z}\nu(u,x).
\]
By Siu's theorem, recalled in \cref{subsec:psh-lelong-levels}, $\nu(u,x)=\nu(u,Z)$ for all $x\in Z$ outside a countable union of proper analytic subsets of $Z$.

For analytic singularities we follow the convention of \cite{SGPT}: we say that a quasi-plurisubharmonic function $u$ has \emph{analytic singularities} if near every point there are holomorphic functions $f_1,\ldots,f_N$, not all identically zero, and a rational number $c>0$ such that
\[
 u=c\log\bigl(|f_1|^2+\cdots+|f_N|^2\bigr)+O(1),
 \qquad c\in\mathbb Q_{>0}.
\]
The remainder is only required to be bounded; if it can be chosen smooth, we say that $u$ has analytic singularities \emph{with smooth remainder}. Given a coherent ideal sheaf $\mathfrak a$ and a common coefficient $c$, we say that $u$ has \emph{analytic type} $(c,\mathfrak a)$ if this holds with local generators $f_1,\ldots,f_N$ of $\mathfrak a$. The choice of generators does not matter: two finite systems of generators of the same ideal change the logarithm only by a bounded amount on smaller relatively compact open sets.

\begin{proposition}\label{prop:analytic-type-calculus}
Analytic singularities are preserved by addition, finite maxima, and multiplication by positive rational numbers. Let $u$ have analytic type $(c,\mathfrak a)$, and let $Y\subset X$ be a connected complex submanifold on which $\mathfrak a$ does not vanish identically. Then the restriction $u|_Y$ is a quasi-plurisubharmonic function on $Y$, in particular not identically $-\infty$, and it has analytic type
\[
 (c,\mathfrak a\mathcal O_Y).
\]
Here $\mathfrak a\mathcal O_Y$ is the image ideal, generated by the restrictions to $Y$ of local sections of $\mathfrak a$.
\end{proposition}
\begin{proof}
For multiplication by a positive rational number, multiply the coefficient. Consider two presentations with coefficients $c_1,c_2$ and ideals $\mathfrak a_1,\mathfrak a_2$, and choose a positive rational number $d$ with $c_1=k d$ and $c_2=\ell d$, where $k,\ell$ are positive integers. Raising a sum of squared moduli of generators of $\mathfrak a_1$ to the power $k$ gives, up to bounded factors, the corresponding sum for $\mathfrak a_1^k$; similarly for $\mathfrak a_2^\ell$. Products of generators of $\mathfrak a_1^k$ and $\mathfrak a_2^\ell$ therefore present the sum with coefficient $d$, since a product of sums of squares is a sum of squares of pairwise products.

\Needspace{7\baselineskip}
For the maximum, apply the elementary inequalities
\[
 \max\{A,B\}\leq A+B\leq2\max\{A,B\},\qquad A,B\geq0,
\]
to the two sums of squared moduli of generators of $\mathfrak a_1^k$ and $\mathfrak a_2^\ell$. Taking logarithms shows that the maximum differs from $d$ times the logarithm of their total by a bounded amount. These are local assertions; on the compact manifold, finitely many charts give the corresponding comparisons of ordinary singularity types.

For the restriction, restrict to $Y$ the two-sided bounds that define the remainder. Since $Y$ is connected and is not contained in the zero set of $\mathfrak a$, the common zero set of the restricted generators is a proper analytic subset of $Y$, so $u|_Y$ is finite on a dense open subset of $Y$. After a smooth local correction, the restriction theorem for plurisubharmonic functions shows that $u|_Y$ is quasi-plurisubharmonic, and the restricted bounds give the asserted presentation.
\end{proof}

For an arbitrary quasi-plurisubharmonic function, a vanishing generic Lelong number along a submanifold does not prevent the restriction from being identically $-\infty$. For example, on a small polydisk in $\mathbb C^n$ with coordinates $z=(z_1,\ldots,z_n)$ and $|z_1|<\mathrm{e}^{-1}$, the function
\[
 u(z)=-\sqrt{-\log|z_1|^2}
\]
is plurisubharmonic, because $t\mapsto-\sqrt{-t}$ is increasing and convex for $t<0$. Its Lelong number vanishes at every point of $\{z_1=0\}$, but its restriction to this hyperplane is identically $-\infty$. For the analytic approximants constructed below, the analytic hypothesis in \cref{prop:analytic-type-calculus} excludes this behavior.

\subsection{Lelong level sets}
\label{subsec:psh-lelong-levels}\index{Siu analyticity theorem}

Siu's analyticity theorem states that, for every quasi-plurisubharmonic function $u$ and every $c>0$, the upper level set $\{x:\nu(u,x)\geq c\}$ is an analytic subset. After a smooth local correction, this is the theorem for closed positive currents, which we use in the form recalled in \cite{SGPT}; see \cref{sec:history-psh} for the original source.

For a closed positive current, the Lelong monotonicity formula also bounds the trace mass of small balls. The resulting Hausdorff estimate shows that a closed positive $(j,j)$-current puts no mass on sets of zero Hausdorff $2(n-j)$-measure.

\Needspace{10\baselineskip}
\begin{corollary}\label{cor:closed-current-hausdorff-null}\index{Closed positive current!Hausdorff support}
Let $T$ be a closed positive $(j,j)$-current on a complex manifold of dimension $n$, with $1\leq j<n$. Its coefficient measures vanish on every Borel set of zero Hausdorff $2(n-j)$-measure.
\end{corollary}
\begin{proof}
In a coordinate chart put $\kappa=\ddc|z|^2$ and $\mu=T\wedge\kappa^{n-j}$. By the Lelong monotonicity formula in bidimension $(n-j,n-j)$, the function $r\mapsto r^{-2(n-j)}\mu(B(x,r))$ is nondecreasing; see \cite[Chapter III, formulas (5.5)--(5.8), pp.~158--159]{DemCADG}. Fix a relatively compact inner chart and a radius $r_0>0$ such that all balls of radius $r_0$ centered in it lie in a fixed compact subset of the chart. Their masses are bounded by the mass of that compact set. Hence
\[
 \mu(B(x,r))\leq Cr^{2(n-j)},\qquad 0<r\leq r_0,
\]
uniformly for centers in the inner chart. A Borel set in the inner chart with zero Hausdorff $2(n-j)$-measure can be covered by balls $B(x_i,r_i)$ with $\sum_ir_i^{2(n-j)}$ arbitrarily small, so it has zero $\mu$-mass. Since $T$ is positive, the total variation of each of its coefficient measures is bounded by a constant multiple of $\mu$. A countable cover by charts proves the assertion.
\end{proof}

The non-strict threshold in Siu's theorem is essential, even on a compact K\"ahler manifold.

\begin{example}\label{ex:strict-lelong-level}\index{Lelong number!strict level sets}
Let $X=\mathbb P^2$, with homogeneous coordinates $[x:y:z]$, and let $\sigma$ be the Fubini--Study form normalized by $\int_X\sigma^2=1$. Put $\|Z\|^2=|x|^2+|y|^2+|z|^2$, $a_j=2^{-j}$ and $\lambda_j=2^{-j-1}$ for $j\geq1$. The global function
\[
 u([Z])=\log\frac{|y|^2}{\|Z\|^2}
       +\sum_{j\geq1}\lambda_j
          \log\frac{|x-a_jz|^2}{(1+a_j^2)\|Z\|^2}
\]
belongs to $\PSH(X,\frac32\sigma)$. For the line $D=\{y=0\}$ and the points $p_j=[a_j:0:1]$, it satisfies
\[
 \{\nu(u,\cdot)\geq1\}=D,
 \qquad \{\nu(u,\cdot)>1\}=\{p_j:j\geq1\}.
\]
The second set is not analytic.
\end{example}
\begin{proof}
For a unit linear form $\ell$, the function $v_\ell=\log(|\ell(Z)|^2/\|Z\|^2)$ is nonpositive and satisfies $\ddc v_\ell=[\ell=0]-\sigma$. Its $L^1(X,\sigma^2)$ norm is finite and independent of $\ell$: the group $U(3)$ acts on homogeneous coordinates, preserves $\sigma$, and acts transitively on unit linear forms. Since $\sum_j\lambda_j=1/2$, the displayed partial sums decrease to an $L^1$ function. Their common curvature lower bound is $-\frac32\sigma$, so their limit is quasi-plurisubharmonic and
\[
 T:=\tfrac32\sigma+\ddc u=[D]+\sum_{j\geq1}\lambda_j[L_j],
 \qquad L_j=\{x-a_jz=0\}.
\]
The series of positive closed currents converges: every line has $\sigma$-trace mass one and the coefficients are summable.

We compute the Lelong numbers from the local mass formula. In affine coordinates $w=(w_1,w_2)$ centered at a point $q$, put $\kappa=\ddc|w|^2$. An affine complex line $H$ at distance $d$ from $q$ satisfies
\[
 \int_{H\cap B(q,r)}\kappa=(r^2-d^2)_+.
\]
Consequently $r^{-2}\int_{H\cap B(q,r)}\kappa$ is at most one and tends to $\mathbf1_H(q)$ as $r\downarrow0$; a line outside the chart contributes zero. Tonelli's theorem and dominated convergence for the summable coefficients therefore give
\[
 \nu(u,q)=\nu(T,q)
       =\mathbf1_D(q)+\sum_{j\geq1}\lambda_j\mathbf1_{L_j}(q).
\]
Here a smooth local correction does not change the Lelong number. This computation includes points where infinitely many lines accumulate.

The lines $L_j$ meet $D$ at the distinct points $p_j$ and meet one another only at $p_\infty=[0:1:0]\notin D$. Thus $\nu(u,p_j)=1+\lambda_j$, whereas $\nu(u,q)=1$ for $q\in D\setminus\{p_j:j\geq1\}$. Outside $D$, the Lelong number is at most $\sum_j\lambda_j=1/2$. The two level-set identities follow. Finally, $p_j\to p_*=[0:0:1]$, and $\nu(u,p_*)=1$. The strict level set is not closed and hence is not analytic.
\end{proof}


\section{The local Bergman estimates}

Demailly's approximation of plurisubharmonic functions by functions with analytic singularities is built from Bergman kernels of weighted spaces of holomorphic functions. In this section we recall the local construction and its estimates. Fix a bounded pseudoconvex domain $\Omega\subset\mathbb C^n$ and a plurisubharmonic function $\psi$ on $\Omega$. Write $\mathrm{d}\lambda$ for Euclidean volume and $B(x,r)$ for the Euclidean ball of radius $r>0$ about $x$. For a positive integer $m$, define the weighted Hilbert space of holomorphic functions
\[
 \mathcal H_m=\left\{f\in\mathcal O(\Omega):
          \|f\|_m^2:=\int_\Omega|f|^2\mathrm{e}^{-m\psi}\,\mathrm{d}\lambda<\infty\right\}.
\]
\Needspace{8\baselineskip}
Choose an orthonormal basis $(g_{m,j})_j$ of $\mathcal H_m$ and define $\beta_m=m^{-1}\log\sum_j|g_{m,j}|^2$. The sum is the squared norm of point evaluation,
\[
 \sum_j|g_{m,j}(x)|^2
   =\sup_{\|f\|_m\leq1}|f(x)|^2,
\]
and is therefore independent of the basis. By the local finite generation part of Bergman regularization, $\beta_m$ has analytic type $(1/m,\mathcal I(m\psi))$ on relatively compact open subsets of $\Omega$. This follows from the generation of the multiplier ideal by weighted $L^2$ holomorphic functions and from local $L^2$ division; see \cite[Proposition 5.7 and Section 14.A]{DemAM11}. We use this local statement, together with the point-extension estimate in the proof below, as a standard input.

\begin{proposition}\label{prop:bergman-local-estimates}
There are constants $C_1,C_2$ independent of $m$ such that, for $B(x,r)\Subset\Omega$,
\begin{equation}\label{eq:bergman-local-estimates}
 \psi(x)-\frac{C_1}{m}
 \leq\beta_m(x)
 \leq\sup_{B(x,r)}\psi+\frac1m\log\frac{C_2}{r^{2n}}.
\end{equation}
Consequently $\beta_m\to\psi$ pointwise and in $L^1_{\mathrm{loc}}$, and
\begin{equation}\label{eq:bergman-lelong-bounds}
 \max\{\nu(\psi,x)-n/m,0\}
 \leq\nu(\beta_m,x)\leq\nu(\psi,x).
\end{equation}
\end{proposition}
\begin{proof}
For $\|f\|_m\leq1$, the submean inequality gives
\[
 |f(x)|^2
 \leq C_2r^{-2n}\int_{B(x,r)}|f|^2\,\mathrm{d}\lambda
 \leq C_2r^{-2n}\mathrm{e}^{m\sup_{B(x,r)}\psi}.
\]
Taking the supremum over the unit ball of $\mathcal H_m$ and then the logarithm proves the upper bound.

For the lower bound, suppose that $\psi(x)>-\infty$. The local Ohsawa--Takegoshi point-extension theorem gives a holomorphic $f$ with $f(x)=1$ and
\[
 \int_\Omega|f|^2\mathrm{e}^{-m\psi}\,\mathrm{d}\lambda
 \leq C \mathrm{e}^{-m\psi(x)},
\]
where the constant is independent of the plurisubharmonic weight and of $m$. After normalizing $f$ in $\mathcal H_m$, we see that the squared norm of point evaluation is at least $C^{-1}\mathrm{e}^{m\psi(x)}$, which proves the lower bound. At a pole of $\psi$ the lower bound holds trivially.

Choosing $r=1/m$ in the upper bound and using upper semicontinuity of $\psi$ proves pointwise convergence. On a fixed relatively compact set, the same estimates bound $\beta_m$ above by a constant and below by $\psi-C_1/m$. Since $\psi$ is locally integrable, dominated convergence also gives convergence in $L^1_{\mathrm{loc}}$.

For the Lelong estimate, take suprema over $B(x,r)$ in the two inequalities, bounding the right-hand side of the upper one by a supremum over $B(x,2r)$. Divide by $\log r^2<0$ and let $r\downarrow0$. After the division, the term $-(2n/m)\log r$ contributes $-n/m$, while the constant terms disappear. This proves the two inequalities in \eqref{eq:bergman-lelong-bounds}; since Lelong numbers are nonnegative, the lower bound may be replaced by its maximum with zero.
\end{proof}

In \cite{DemAM11}, the local weight is $\mathrm{e}^{-2m\phi}$ and Lelong numbers are normalized by $\log|z|$. With $\psi=2\phi$, these conventions give exactly the formulas of \cref{prop:bergman-local-estimates}; in particular, the loss is $n/m$ for the weight $\mathrm{e}^{-m\psi}$ and the normalization by $\log|z|^2$.

The same construction and estimates hold for every real $m>0$. The coefficient $1/m$ of the resulting logarithmic presentation is rational exactly when $m$ is rational, and we speak of an analytic type only in that case. The global sequence constructed below is indexed by integers.

For a quasi-plurisubharmonic function $u$, the generators of $\mathcal I(mu)$ also control the multiplier ideals of $u$ at nearby weights; a small loss in the weight compensates for the finite Bergman index.

\begin{lemma}\label{lem:bergman-controlled-ideals}\index{Multiplier ideal!controlled approximation}
Let $u$ be quasi-plurisubharmonic, let $m>0$, and let $a_m$ be a function which locally satisfies $a_m=m^{-1}\log\sum_k|g_k|^2+O(1)$, where the $g_k$ generate $\mathcal I(mu)$. For $t>0$ and $s>1$,
\[
 \mathcal I(sta_m)\subseteq\mathcal I(tu)
 \qquad\text{whenever}\qquad m\geq\frac{st}{s-1}.
\]
Moreover, $u\preceq a_m$ locally, and hence $\mathcal I(tu)\subseteq\mathcal I(ta_m)$ for every $t>0$.
\end{lemma}
\begin{proof}
On a sufficiently small coordinate neighborhood, choose generators $g_1,\ldots,g_N$ of $\mathcal I(mu)$ and put $B=(\sum_k|g_k|^2)\mathrm{e}^{-mu}$. After shrinking the neighborhood, $B$ is integrable and $a_m=m^{-1}\log\sum_k|g_k|^2+O(1)$ with a bounded remainder. If $f\in\mathcal I(sta_m)$, shrink once more so that $f$ is bounded and $A=|f|^2\mathrm{e}^{-sta_m}$ is integrable. Then, almost everywhere,
\[
 |f|^2\mathrm{e}^{-tu}
 \leq C|f|^{2-2/s}A^{1/s}B^{t/m}.
\]
H\"older's inequality with exponents $s$ and $s/(s-1)$ applies because $st/(m(s-1))\leq1$, and it shows that the left side is locally integrable, also in the case of equality in the bound for $m$. Finally, the local Bergman lower estimate of \cref{prop:bergman-local-estimates} gives $u\preceq a_m$ locally; bounded smooth corrections affect neither this comparison nor any of the ideals.
\end{proof}

\Needspace{23\baselineskip}
\section{The global curvature bound}

We return to the compact manifold $X$ with its Hermitian form $\omega$. Demailly's global regularization theorem glues the local Bergman constructions while controlling the curvature error. We use it as a classical input in the following form.

\begin{theorem}[Demailly regularization]\label{thm:demailly-hermitian-input}\index{Regularization!analytic singularities}
Let $u$ be quasi-plurisubharmonic on $X$ and let $\gamma$ be a continuous real $(1,1)$-form with $\ddc u\geq\gamma$. There are quasi-plurisubharmonic functions $u_m$, defined for all sufficiently large positive integers $m$, and numbers $\epsilon_m\downarrow0$ such that
\begin{enumerate}
\item $u_m$ has analytic type $(1/m,\mathcal I(mu))$;
\item $u_m\geq u$ and $u_m\to u$ pointwise and in $L^1(X)$;
\item $\ddc u_m\geq\gamma-\epsilon_m\omega$;
\item $\max\{\nu(u,x)-n/m,0\}\leq\nu(u_m,x)\leq\nu(u,x)$ for all $x$.
\end{enumerate}
The sequence $(u_m)$ need not be monotone.
\end{theorem}

This is the form of \cite[Theorem 14.12]{DemAM11} used in this book. The local generators in the construction are the weighted Bergman generators of $\mathcal I(mu)$; smooth local corrections do not change this ideal. On overlaps the local Bergman sums are comparable, so the glued logarithm keeps the analytic type $(1/m,\mathcal I(mu))$. The construction freezes the continuous lower bound $\gamma$ on sufficiently small coordinate balls, and its curvature error combines the modulus of continuity of $\gamma$ with an overlap term that tends to zero as $m$ increases. It is this global curvature estimate, and not merely the local convergence of \cref{prop:bergman-local-estimates}, that we take from the theorem.

\begin{corollary}\label{cor:nonclosed-analytic-approximation}
Let $\theta$ be a smooth real $(1,1)$-form on $X$ and $u\in\PSH(X,\theta)$. The approximants $u_m$ of \cref{thm:demailly-hermitian-input} for $\gamma=-\theta$ satisfy
\[
 \theta+\ddc u_m\geq-\epsilon_m\omega.
\]
In particular, they have analytic type $(1/m,\mathcal I(mu))$.
\end{corollary}
\begin{proof}
Apply \cref{thm:demailly-hermitian-input} to $u$ with the lower bound $\gamma=-\theta$; the theorem does not require $\gamma$ to be closed. Adding $\theta$ to the curvature inequality gives the assertion. In the normalization of \cite{DemAM11}, the theorem is applied to the potential $u/2$, since $\frac{\mathrm{i}}{\pi}\partial\bar\partial(u/2)=\ddc u$.
\end{proof}

Thus the approximants keep the curvature lower bound up to an error that tends to zero. Their multiplier ideals determine the limit of their ordinary envelopes; see \cref{thm:dyadic-analytic-envelope-limit}.


\subsection{Regularization of closed currents}
\label{subsec:psh-closed-regularization}

We also use Demailly's regularization of closed currents, in a smooth and in an analytic form; the sources are given in \cref{sec:history-psh}. Let $T=\eta+\ddc v$ be a closed real $(1,1)$-current, where $\eta$ is a smooth closed form and $v$ is quasi-plurisubharmonic, and set $\nu(T,x)=\nu(v,x)$. Suppose that $T\geq\gamma$ for a smooth real $(1,1)$-form $\gamma$. Smooth regularization gives smooth representatives $T_j=\eta+\ddc v_j$ with
\[
 T_j\geq\gamma-C\lambda_j\omega,
\]
where $C>0$ is a fixed constant and the continuous functions $\lambda_j$ decrease pointwise to $\nu(T,\cdot)$. If all Lelong numbers of $T$ vanish, Dini's theorem shows that $\lambda_j\to0$ uniformly, so the loss of positivity becomes uniformly small. Analytic regularization \cite[Corollary~14.13]{DemAM11} instead gives representatives $S_j$ with analytic singularities and
\[
 S_j\geq\gamma-\varepsilon_j\omega,\qquad
 \varepsilon_j\downarrow0,\qquad
 \nu(S_j,x)\leq\nu(T,x).
\]
Thus analytic regularization preserves strict positivity up to a small loss and creates no positive Lelong number at a point where the Lelong number of $T$ vanishes.

\subsection{Smooth remainders after resolution}

For arguments on a modification, the bounded remainder in an analytic presentation is sometimes insufficient. The following version of regularization gives a smooth remainder after resolution.

\Needspace{13\baselineskip}
\begin{theorem}[Demailly regularization]\label{thm:resolvable-regularization-input}
Let $\alpha$ be a smooth real $(1,1)$-form and let $u\in\PSH(X,\alpha)$. There are numbers $\delta_j\downarrow0$ and functions $\widetilde u_j\in\PSH(X,\alpha+\delta_j\omega)$ decreasing to $u$ with the following properties. Each $\widetilde u_j$ is smooth outside an analytic pole set, and there is a modification $\pi_j\colon Y_j\to X$, obtained by a finite sequence of blowups with smooth centers, such that locally on $Y_j$
\[
 \pi_j^*\widetilde u_j=\sum_\ell a_\ell\log|s_\ell|^2+b,
 \qquad a_\ell>0,\quad b\in C^\infty.
\]
Here the $s_\ell$ are local holomorphic defining functions of divisors, the $a_\ell>0$ are real numbers, and $b$ is a smooth real function. The divisors may be taken to have simple normal crossings.
\end{theorem}

This classical theorem is applied to the closed current $\ddc u$ with the smooth lower bound $-\alpha$; see \cref{sec:history-psh} for its source. The lower bound need not be closed, and $X$ need not be K\"ahler. What matters here is the smooth remainder $b$ after resolution; an analytic type, whose remainder is merely bounded, would not suffice. The sequence $(\widetilde u_j)$ need not be the multiplier-ideal sequence of \cref{thm:demailly-hermitian-input}, and we do not claim that the conclusions of the two theorems hold for one and the same sequence. If $\theta+\ddc u\geq\varepsilon\omega$ for some $\varepsilon>0$, applying the theorem with $\alpha=\theta-\varepsilon\omega$ and discarding the terms with $\delta_j\geq\varepsilon$ gives approximants with $\theta+\ddc\widetilde u_j\geq(\varepsilon-\delta_j)\omega>0$; thus strict positivity is preserved.
